\chapter{Conclusion and Future Work}
\label{ch:conclusion}

\section{Summary of Contributions}
\label{sec:conclusion_summary}

This thesis presented a privacy-preserving, federated anomaly detection
framework for distributed microservice architectures. The framework
integrates six core components, each addressing a distinct challenge in
operational monitoring:

\begin{enumerate}
    \item \textbf{Federated LSTM-Autoencoder for Anomaly Detection.} A
    two-layer LSTM-based autoencoder (Chapter~\ref{ch:model_architecture})
    trained across distributed edge nodes using Federated Averaging
    (Chapter~\ref{ch:federated_learning}). The model learns normal system
    behavior from real-world telemetry (SMD dataset) and detects anomalies
    via reconstruction error, achieving an F1-score of 0.839 with an
    AUC-ROC of 0.950, significantly outperforming a centralized baseline
    (Section~\ref{subsec:eval_detection_results}).

    \item \textbf{Differential Privacy Integration.} Client-level gradient
    clipping and Gaussian noise injection provide formal DP guarantees for
    model updates shared during federated training. Experimental evaluation
    demonstrates the privacy--accuracy trade-off, with detection AUC-ROC
    decreasing from 0.972 (no DP) to 0.431 at the strictest tested privacy
    budget ($\varepsilon \approx 0.1$), establishing baseline behavior for DP-SGD
    in distributed telemetry environments
    (Section~\ref{subsec:eval_dp_impact}).

    \item \textbf{Causal-Graph-Based Root Cause Analysis.} A causal dependency
    graph, constructed from distributed traces
    (Chapter~\ref{ch:root_cause_analysis}), enables automated root cause
    identification when anomalies are detected. The PageRank-inspired scoring
    algorithm achieves a Top-1 accuracy of 42\%, a Top-3 accuracy of 89\%,
    and an MRR of 0.653, outperforming a random baseline (35\% Top-1, 0.601
    MRR) on the Train-Ticket microservice topology
    (Section~\ref{subsec:eval_rca_results}).

    \item \textbf{Q-Learning Adaptive Thresholding.} A per-service tabular
    Q-learning agent (Chapter~\ref{ch:adaptive_thresholding}) dynamically
    adjusts the anomaly detection threshold based on detection outcomes and
    SLO compliance feedback, reducing reliance on manually configured static
    thresholds. The agent's threshold trajectory demonstrates adaptation to
    sinusoidal traffic patterns over a 24-hour simulation, with
    $\varepsilon$-greedy exploration decaying from 0.10 to 0.05
    (Section~\ref{subsec:eval_threshold_convergence}).

    \item \textbf{Efficient FL Communication Stack.} An adaptive serialization
    and compression pipeline (Chapter~\ref{ch:communication}) reduces model
    update payloads by up to 61\% using Zstd compression on JSON/Protobuf
    payloads. For binary formats, the controller switches to LZ4 under high CPU load,
    providing $5.8\times$ faster compression while sacrificing minimal
    payload reduction (Section~\ref{subsec:eval_adaptive_compression}).

    \item \textbf{Fault-Tolerant Edge Design.} Two node profiles---lightweight
    (stateless) and standard (persistent SQLite cache)---enable crash recovery
    and seamless FL participation across heterogeneous edge deployments
    (Chapters~\ref{ch:federated_learning} and~\ref{ch:reliability}).
\end{enumerate}

\section{Conclusion}
\label{sec:conclusion_conclusion}

The experimental evaluation (Chapter~\ref{ch:evaluation}) demonstrates that
the proposed framework achieves the following key results:

\begin{itemize}
    \item \textbf{High accuracy on real-world data.} The federated LSTM-AE
    achieves an F1-score of 0.839 on the Server Machine Dataset, surpassing
    the centralized baseline (0.496). This confirms that the federated
    aggregation of distributed representations provides stronger generalization
    than single-node centralization on realistic telemetry.
    The evaluation also quantifies the DP trade-off, showing AUC-ROC gracefully
    degrading under strict privacy budgets.

    \item \textbf{Efficient communication.} Binary serialization (TorchSave)
    produces payloads $5.4\times$ smaller than JSON or Protobuf. Zstd
    compression achieves an additional 9\% reduction on binary payloads and
    up to 61\% on text-based formats. The full FL communication
    pipeline---compression, transmission, and aggregation---accounts for less
    than 3\% of total round time, with local training dominating at over 97\%.

    \item \textbf{Actionable root cause identification.} The RCA engine
    identifies the true root cause in its top-3 ranking for 89\% of
    fault-injection events on Train-Ticket, providing operators with a focused
    investigation set rather than an undifferentiated list of anomalous services.
    The framework consistently outperforms a random baseline in MRR (0.653 vs.\
    0.601).

    \item \textbf{Linear scalability.} FedAvg aggregation scales linearly
    from 5.1\,ms at 5~nodes to 990.1\,ms at 1\,000~nodes, with bounded memory
    consumption ($\approx 42$\,MB). This empirical scalability confirms the
    coordinator is not a bottleneck up to at least 1\,000 concurrent
    participants.
\end{itemize}

These results indicate that federated learning, differential privacy, and
causal RCA can be combined into a practical, deployable system for anomaly
detection in distributed microservice environments without requiring raw data
centralization or manual threshold configuration.

\section{Limitations}
\label{sec:conclusion_limitations}

Despite the positive results, several limitations should be acknowledged:

\begin{itemize}
    \item \textbf{Cold-start problem.} The Q-learning threshold agent requires
    an initial exploration phase during which alert quality may be suboptimal.
    The benchmark results show negative cumulative reward during the first
    12~hours, indicating that a warm-up period is necessary before the policy
    stabilizes. Newly onboarded services experience this cold-start phase
    independently.

    \item \textbf{IID data assumption.} FedAvg assumes that local datasets are
    approximately IID across clients
    \cite{li2020federated, mcmahan2017communication}. In production
    deployments where edge nodes monitor services with substantially different
    workload characteristics, the aggregated global model may underperform
    compared to locally specialized models.

    \item \textbf{RCA accuracy on deep dependency chains.} The upstream
    analysis approach performs well when the root cause has few anomalous
    predecessors but struggles with services like \texttt{db-service} that sit
    at the bottom of many dependency chains. The Top-1 accuracy of 42\% on
    Train-Ticket indicates room for more sophisticated causal inference methods.

    \item \textbf{DP noise and convergence.} Although detection AUC-ROC is
    acceptable under moderate DP, the training loss plateaus at $\approx 0.0014$
    compared to 0.0000 without DP, and extreme privacy budgets ($\varepsilon \approx 0.1$)
    degrade AUC-ROC to 0.431. For larger models or more complex anomaly patterns, this
    convergence gap translates into reduced detection sensitivity.

    \item \textbf{Single-coordinator architecture.} The current design relies
    on a single coordinator instance, which, despite crash-recovery mechanisms
    (Chapter~\ref{ch:reliability}), remains a potential availability bottleneck
    under extreme node counts.
\end{itemize}

\section{Future Work}
\label{sec:conclusion_future}

\subsection{Transfer Learning}
\label{subsec:future_transfer_learning}

The current framework trains the LSTM-AE from scratch for each deployment.
A promising direction is to pre-train the model on a large, generic corpus of
system metric data---such as the Alibaba Cluster Trace
dataset \cite{zhang2019tracing}---and then fine-tune on local edge data using
few-shot adaptation. This would reduce the number of FL rounds required for
convergence and improve detection quality for newly onboarded services that
have limited local training data.

\subsection{Hierarchical Federated Learning}
\label{subsec:future_hierarchical_fl}

For deployments exceeding 1\,000~edge nodes, a two-tier hierarchical FL
architecture would reduce the communication and computational load on the
global coordinator. In this design:

\begin{itemize}
    \item \textbf{Cluster-level sub-coordinators} aggregate model updates from
    nearby edge nodes within a geographic or administrative boundary.
    \item \textbf{The global coordinator} aggregates from sub-coordinators
    only, reducing the number of incoming connections by a factor proportional
    to the number of clusters.
\end{itemize}

This approach is compatible with the existing FedAvg aggregation logic and
requires only routing changes rather than algorithmic modifications
\cite{kairouz2021advances}.

\subsection{Enhanced Privacy Mechanisms}
\label{subsec:future_privacy}

The current implementation uses central DP, where the coordinator is assumed
to be honest-but-curious. Three complementary privacy mechanisms could
strengthen the threat model:

\begin{itemize}
    \item \textbf{Secure Multi-Party Computation (SMPC)}: Enables aggregation
    without revealing individual updates to the coordinator, removing the
    trusted-coordinator assumption \cite{bonawitz2019towards}.

    \item \textbf{Homomorphic Encryption}: Allows the coordinator to aggregate
    encrypted gradient updates without decryption, providing strong
    confidentiality guarantees at the cost of increased computational overhead.

    \item \textbf{Local Differential Privacy (LDP)}: Applies noise at the
    data collection level rather than at the gradient level, providing
    stronger per-sample privacy but typically requiring larger datasets for
    equivalent model utility \cite{dwork2014algorithmic}.
\end{itemize}

\subsection{Improved Root Cause Analysis}
\label{subsec:future_rca}

Several enhancements could improve RCA accuracy beyond the current Top-1 rate
of 42\%:

\begin{itemize}
    \item \textbf{Multi-modal causal graph}: Incorporating structured log
    events and deployment metadata alongside trace data would provide richer
    causal signals. Error logs, memory allocation events, and container
    restarts could be modeled as additional edge types in the causal graph.

    \item \textbf{Temporal causal graph}: Annotating edges with propagation
    delays (derived from span durations) would enable the scoring algorithm
    to reason about temporal precedence, favoring services whose anomalies
    consistently precede downstream anomalies \cite{fu2020root}.

    \item \textbf{LLM-assisted incident analysis}: Large language models could
    be used to generate human-readable incident summaries from the RCA output,
    correlating root-cause rankings with relevant log snippets and suggesting
    remediation actions based on historical resolution patterns.
\end{itemize}

\subsection{Deep Q-Network for Adaptive Thresholding}
\label{subsec:future_dqn}

The current tabular Q-learning implementation discretizes the state space into
coarse buckets, limiting the agent's ability to distinguish fine-grained
detection regimes. Replacing the Q-table with a Deep Q-Network (DQN)
\cite{turn0search6} would enable:

\begin{itemize}
    \item Continuous state representation without manual discretization.
    \item Richer state features (e.g., rolling reconstruction error
    statistics, request volume, time-of-day encoding).
    \item Potentially faster convergence through experience replay and target
    network stabilization.
\end{itemize}

This enhancement would be particularly valuable for services with complex,
multi-modal traffic patterns where tabular discretization loses important
signal.

\subsection{Multi-Modal Anomaly Detection}
\label{subsec:future_multimodal}

The current LSTM-AE operates exclusively on numerical metric time series.
Modern microservice deployments produce multiple telemetry modalities---metrics,
structured logs, and distributed traces---each capturing different aspects of
system behavior. Extending the LSTM-AE to jointly model these modalities
through a multi-modal encoder (e.g., separate specialized encoders for each
modality with a shared latent space) could improve detection accuracy for
complex anomalies that manifest across multiple telemetry types simultaneously.
Recent datasets such as AnoMod \cite{zhang2019tracing} provide five-modality
benchmarks suitable for evaluating such architectures.
